\section*{Chapter \ref{ch:dv:barriers-and-digitals} --- Barriers and Digitals}

\begin{solution}{exo:dv:barriers-and-digitals:1}
Exactly one of the two pays 1 at expiry, whatever happens, so together they pay 1 for sure: a zero-coupon bond,
$e^{-rT}$. Their deltas therefore add up to zero: the digital put's delta is minus the call's.
\end{solution}

\begin{solution}{exo:dv:barriers-and-digitals:2}
The digital is $-\partial C/\partial K$. With a downside skew the call's volatility falls as the strike rises, so the call
loses value with the strike faster than at a constant volatility, and $-\partial C/\partial K$ gains the term
$-\mathcal V\,\partial\sigma/\partial K>0$. In terms of probabilities: the skewed density has a fat left tail and, to keep
the forward fixed, more mass just above the forward, so the probability of finishing above an at-the-money strike is higher
than the flat model's.
\end{solution}

\begin{solution}{exo:dv:barriers-and-digitals:3}
In-out parity is violated: $7.23+1.70=8.93\ne8.83$. With zero rebates the two must add up to the vanilla; the knock-in
should be 1.60, or one of the prices uses different inputs (barrier, monitoring or volatility).
\end{solution}

\begin{solution}{exo:dv:barriers-and-digitals:4}
$K/H=1.11$ puts struck at $H^2/K=81$ are worth 1.498, and the down-and-in call is worth 1.498. If the barrier is never hit
the index stays above 90 and ends above 90, hence above the puts' strike of 81: the puts expire worthless, like the
knock-in.
\end{solution}

\begin{solution}{exo:dv:barriers-and-digitals:5}
The factor is $e^{0.5826\times0.2\times\sqrt{1/252}}=e^{0.0073}$: the barrier at 90 moves down to 89.34 and the barrier at 120
up to 120.88, away from the spot in both cases.
\end{solution}

\begin{solution}{exo:dv:barriers-and-digitals:6}
As the spot rises towards 120 the call's intrinsic value grows, but so does the probability of touching the barrier and
losing everything; near the barrier the second effect dominates, so the value falls with the spot and the delta is
negative. The value peaks where the two effects balance, well below the barrier (near 111 with one month left, where it is worth 7.68).
\end{solution}

\begin{solution}{exo:dv:barriers-and-digitals:7}
Continuous 0.162. Daily: Monte Carlo 0.205, shifted 0.207; weekly 0.262 and 0.270; monthly 0.383 and 0.430; four dates 0.564
and 0.735. The shift works daily, and degrades faster than for the call as monitoring thins out, because the put's
payoff is large at the barrier ($K-H=10$), where the expansion is least accurate.
\end{solution}

\begin{solution}{exo:dv:barriers-and-digitals:8}
A stop-loss order at the barrier is filled at the next available price. If the market gaps through the level (an
announcement, a policy change, an overnight move), the fill is far beyond it and the hedge's P\&L is the gap times the
position, as on 15 January 2015. \omterm{def:dv:barriers-and-digitals:gap}{Gap risk} must be priced (jump models) and limited (by exposure at each barrier level),
not assumed away.
\end{solution}

\begin{solution}{pb:dv:barriers-and-digitals:1}
\textbf{1.} 4.84 million at 16.4\%; 5.77 million with the smile.
\textbf{2.} 0.93 million: $-\mathcal V\,\partial\sigma/\partial K$, positive because the skew makes volatility fall with the
strike.
\textbf{3.} 1.25 million per index point, equivalent to 6.25 billion of index.
\textbf{4.} It would have to be bought or sold within a minute and reversed at every tick through the strike: far beyond
the market's depth, and the hedge would lose the spread each time.
\textbf{5.} For the holder, large and positive just below the strike and large and negative just above it; the seller has
the opposite.
\textbf{6.} Long calls at $K-\varepsilon$, short calls at $K$, in size $N/\varepsilon$: it pays at least the digital
everywhere, so the seller who holds it is never short at expiry.
\textbf{7.} $\varepsilon=10\,000\,000/50\,000=200$ points: the spread's payoff changes by at most 50\,000 per point.
\textbf{8.} 6.66 million; the overhedge is 0.89 million over the digital's 5.77 million.
\textbf{9.} 0.25, 0.48 and 1.56 million.
\textbf{10.} 5 million, half the digital's notional, instead of nothing: the overhedge pays for this region.
\textbf{11.} 7.23 and 1.60.
\textbf{12.} 7.48 (standard error 0.03) by simulation and 7.45 by the shift.
\textbf{13.} 6.15 against 6.51: \omterm{def:dv:local-volatility:model}{local volatility} rises as the index falls towards the barrier, so the barrier is hit more
often than at a flat volatility.
\textbf{14.} 0.141.
\textbf{15.} 0.592 at the hit, 0.581 at expiry: paying at the hit saves the discounting from the hit to expiry.
\textbf{16.} The client buys a payoff that dominates the digital; the extra price pays for the extra payoff between
4\,800 and 5\,000 and for a hedge the desk can actually run.
\textbf{17.} When the desk can accept some pin risk in exchange for a price close to the digital's, for instance when it
can offset it against other positions struck nearby; a centred spread is cheaper but can leave the seller short at
expiry.
\textbf{18.} Statically where possible (a calendar of options at the barrier), with a limit on the delta at the barrier and a
reserve for gaps; dynamic hedging alone fails exactly when the barrier is crossed by a jump.
\textbf{19.} 200 index points: a call spread of 50\,000 per point from 4\,800 to 5\,000, which charges the client an
overhedge of 0.89 million on the 10-million digital (5.77 million).
\textbf{20.} A discontinuity, in the payoff or at the barrier, where the delta becomes unbounded and the hedge must be
replaced by an overhedge, a static portfolio or a reserve.
\end{solution}

\begin{solution}{iq:dv:barriers-and-digitals:1}
Price it as minus the strike derivative of the call price, which brings in the smile's slope. Hedge it as a call spread:
book the sold digital at the price of the dominating spread (the overhedge) and delta-hedge the spread, whose delta is
bounded.
\iqlookfor{the smile term and the overhedge.}
\end{solution}

\begin{solution}{iq:dv:barriers-and-digitals:2}
The killed density of the log-price at a lower barrier is the free density minus its reflection in the barrier, weighted
by the drift factor. Integrating the call payoff against it gives
$C_{\mathrm{BS}}(S)-(H/S)^{2\nu/\sigma^2}C_{\mathrm{BS}}(H^2/S)$ for $K\ge H$; check that it solves the \omterm{def:dv:black-scholes-three-ways:pde}{Black--Scholes equation}, is zero
on the barrier and has the right payoff.
\iqlookfor{the image, and the boundary check.}
\end{solution}

\begin{solution}{iq:dv:barriers-and-digitals:3}
A digital is $-\partial C/\partial K$, so it depends on the smile's slope at the strike: $-\mathcal V\partial\sigma/\partial K$.
With a downside skew that term is positive for a digital call (worth more than its Black--Scholes price) and negative for a
digital put.
\iqlookfor{the derivative and the sign.}
\end{solution}

\begin{solution}{iq:dv:barriers-and-digitals:4}
No: a daily-monitored knock-out misses crossings between observations, so it is worth more than the continuous formula. Compare
with the formula at the barrier shifted by $e^{\beta\sigma\sqrt{\Delta t}}$ away from the spot; if the two agree within the
simulation error, the Monte Carlo is right.
\iqlookfor{discrete against continuous monitoring, and the shift.}
\end{solution}

\begin{solution}{iq:dv:barriers-and-digitals:5}
The loss if the underlying jumps through a barrier or strike where the position's value or hedge changes abruptly. Measure it
by scenario: the book's P\&L for jumps of several sizes across each barrier level, with hedges assumed filled at the post-jump
price; limit the exposure concentrated at any one level, and price it with jump models or reserves.
\iqlookfor{the scenario, concentration at levels, and pricing.}
\end{solution}

\begin{solution}{iq:dv:barriers-and-digitals:6}
With zero carry and a symmetric smile, hold $K/H$ puts struck at $H^2/K$: on the barrier they are worth exactly the call, so
at the hit swap them for it; if never hit they expire worthless. It fails with carry (the symmetry needs a driftless forward),
with a skew (the puts below the barrier are priced at a higher volatility than the call above it), and with gaps through the
barrier.
\iqlookfor{the symmetry, and its three failure modes.}
\end{solution}
